\subsection{The unscaled radial envelope: a growing local range and an exact lifting failure}
\label{arch136:radial-range}

\paragraph{ARCH-136-RADIAL-RANGE.}
Let \(p\) be prime, \(R=\mathbb Z_{(p)}\), and retain the graded
free module \(V\), the disk algebra \(D\), and the rational bar map
\(\rho\) of ARCH-136-SCALED-BAR. Here no coefficient scaling is
performed. Define
\[
 \widetilde E_R(V)=
 \operatorname{im}\bigl(S_R\otimes_R\mathsf B(\Gamma_R(V))
                  \xrightarrow{\rho}S_{\mathbb Q}\otimes D_{\mathbb Q}\bigr).
 \tag{RR.1}
\]
Then \(\widetilde E_R(V)\) is a strict differential PD
\(S_R\)-subalgebra, natural for the same linear maps as before,
and
\[
 S_R\otimes D_R\ \subseteq\ \widetilde E_R(V),\qquad
 H^*(\widetilde E_R(V))=S_R[0].
 \tag{RR.2}
\]
Moreover its underlying form-degree components satisfy
\[
 \widetilde E_R(V)^q=(S_R\otimes D_R)^q
                         \quad(0\leq q<p).
 \tag{RR.3}
\]
For rank at least \(p\) in form degree two, the bound \(q<p\)
is sharp for this particular envelope.

\paragraph{Proof.}
The proof of multiplicativity and PD preservation of \(\rho\)
in ARCH-136-SCALED-BAR did not require coefficient scaling.
Its image is therefore a differential PD subalgebra. It contains
\(b_v=\rho(b_v)\) and \(a_v=\rho([b_v])\), hence contains the
free PD algebra they generate. The insertion identity (SB.6)
shows that \(H=h/r\) preserves this unscaled image, with no
integrality assertion about \(H\) on the larger ambient algebra.
It is integral on the image because the inserted bar word is
defined over \(R\). The Euler identity (SB.1) gives the
\(S_R\)-linear contraction, proving (RR.2).

A bar word of total arity \(r\) and length \(k\) has form degree
at least \(2r-k\geq r\). Its denominator in (SB.2) divides
\(r!\). Thus every word of form degree \(q<p\) has a denominator
which is a unit in \(R\), and its image already lies in
\(S_R\otimes D_R\). Together with the first inclusion in (RR.2),
this proves (RR.3). The argument remains valid with arbitrarily
many generators and arbitrary positive generator weights.

For distinct form-two generators \(b_1,\ldots,b_p\), one has
\[
 \rho([b_1|\cdots|b_p])=\frac{a_1\cdots a_p}{p!}.
 \tag{RR.4}
\]
The numerator is an exterior basis element. It is not divisible
by \(p\) in \(D_R\), whereas (RR.4) is an element of form
degree \(p\) of the envelope. This proves the sharpness assertion.
\(\square\)

\paragraph{ARCH-136-RADIAL-TARGET.}
Even for the standard admissible theory \(B_R\), and even when
the image of every \(b_v\) is zero, a map from a guarded free
horizontal source need not extend across its unscaled radial
envelope. Thus (RR.3) supplies a growing degree range but does
not supply the target-compatible local completion required in
the integral common-source problem.

\paragraph{Proof.}
Fix \(m\geq2\). In simplicial degree \(p\), write the standard
theory as
\[
 (B_R)_p=S_R\langle x_1,\ldots,x_p\rangle_{\rm PD}
                    \otimes\Lambda_R(e_1,\ldots,e_p),
 \qquad dx_i=e_i,
 \tag{RR.5}
\]
where all displayed generators have weight one, the \(x_i\)
have form degree zero, and the \(e_i\) have form degree one.
The omitted coordinate is \(x_0=T-\sum_{i=1}^p x_i\).
Consider the closed form-one elements of weight \(m\)
\[
                       \alpha_i=x_i^{[m-1]}e_i
                         \quad(1\leq i\leq p).
 \tag{RR.6}
\]
For the rank-\(p\) disk algebra, the assignments
\(a_i\mapsto\alpha_i\), \(b_i\mapsto0\) define a strict
differential PD map. If it extended to its radial envelope, the
image \(y\) of (RR.4) would satisfy
\[
             p!y=\prod_{i=1}^p x_i^{[m-1]}e_1\cdots e_p.
 \tag{RR.7}
\]
The right side is one of the free \(S_R\)-module basis
monomials in (RR.5), with coefficient one. Its reduction modulo
\(p\) is nonzero. The left side reduces to zero. This is a
contradiction.

These data can occur in an actual guarded pure horizontal
construction, rather than only in an isolated disk algebra.
Let
\[
 M=\mathbb Z[\Delta[p]]/\mathbb Z[V_p],
 \tag{RR.8}
\]
where \(V_p\) is the simplicial subset of vertices. Give each
of \(p\) copies of \(M\) a lower vertical generator in form
degree one and weight \(m\), and its derivative in form degree
two. The resulting free differential PD extension of \(S\)
is built by attaching the nondegenerate simplices of \(M\)
in increasing dimension. Each is exactly a horizontal sphere
to disk attachment with the vertical unit disk \(E_1\).
There are no new generators in simplicial degree zero. Since
the standard form-one row is acyclic for weights at least two,
the chosen weight is guarded.

Each \(\alpha_i\in B_p^{m,1}\) defines a simplicial module map
\(\mathbb Z[\Delta[p]]\to B_*^{m,1}\). It vanishes on every
vertex, because positive forms in simplex zero are zero, so
it factors through (RR.8). The derivative partner is sent to
zero because \(d\alpha_i=0\). These assignments define the
strict simplicial map of this initial free source segment to
\(B\). Continuing the guarded horizontal attachments gives a
pure horizontal source \(H\to B\) of ARCH-136-DISK-SUPPORT.
At simplicial degree \(p\), the top representable generator
in each copy is an independent basis generator and has image
\(\alpha_i\). Equation (RR.7) still forbids any extension that
adjoins its radial-envelope element.

In contrast, the one-variable native Moore towers have no such
lifting obstruction for these particular images: all their
higher \(c_r,t_r\) may be sent to zero, because every \(b_i\)
maps to zero. Therefore the obstruction in (RR.7) is the extra
mixed divisibility imposed by this radial envelope. It is not
a counterexample to arbitrary forward completion, nor to the
existence of a compatible local closure whose new form degrees
tend to infinity with \(p\). \(\square\)

\paragraph{Current consequence for the prime-degree route.}
Pure horizontal sources already have no \(p\)-local positive
vertical cohomology below \(2p\). Formula (RR.3) shows that
linear coherence and a growing degree range can coexist in a
natural acyclic disk envelope. Formula (RR.7) shows exactly why
that positive construction cannot be substituted for a repair
retaining maps to the target theories. A successful local
repair must retain the homotopies witnessing divisibility,
as the Moore pairs do, instead of imposing all the radial
divisibilities as strict algebraic relations. No preservation
theorem for those coherences is asserted here.
